\scanpage{61}

Le résultat qui suit, qui résoud le``problème du cycle invariant'' dans le cas du $H^1$, est essentiellement une reformulation d'un résultat de M.RAYNAUD [\quad]:
\margin{$\times$}

\uline{Théorème} / Soient $S$ un trait strictement local, $f:X\unsure{\mathbin{-}|\longrightarrow}S$ un schéma \deleted{propre}\inserted{projectif} sur $S$, régulier, tel que l'on ait $\underline{O}_S\xrightarrow{\sim}f_*(\underline{O}_X)$, (ce qui implique que $X$ est plat sur $S$, et que $X_0$ est un diviseur sur $X$), \deleted{\unsure{soient $(D_i)_{i\in J}$ les}}\inserted{$(D_i)_{i\in J}$ la famille des} composantes irréductibles réduites de $X_0$, considérées comme diviseurs sur $X$, et posons \unsure{$\times$}
\[
(1)\qquad X_0=\sum_{i\in J}d_iD_i,\qquad\text{où }d_i\geq1,
\]
\[
(2)\qquad d=\operatorname{pgcd}(d_i)_{i\in J}.
\]
Soient
\[
(3)\qquad\Delta\simeq\underline{Z}^{J}
\]
la catégorie des diviseurs $D$ sur $X$ tels que $\operatorname{supp}D\subset X_0$, \deleted{et}
\[
(4)\qquad\deleted{\unsure{Ixx\underline Z}}\ \Gamma\simeq/\underline Z\cdot X_0\hookrightarrow\operatorname{Pic}(X)
\]
le \deleted{\unsure{cat}} sous-groupe de $\operatorname{Pic}(X)$ image de $\Delta$ par l'homomorphisme $D\mapsto\operatorname{cl}(\underline{O}_X(D))$ de $\Delta$ dans $\operatorname{Pic}(X)$, et considérons l'homomorphisme \deleted{composé}
\[
(5)\qquad w:\Gamma\longrightarrow\operatorname{NS}(X_0)\overset{\mathrm{dfn}}{=}\operatorname{Im}(\operatorname{Pic}(X_0)\longrightarrow\underline{\operatorname{NS}}_{X_0/k}(k)),
\]
composé des homomorphismes évidents
\diagram{\begin{tikzcd}
\Gamma\arrow[r,hook,"i"]&\operatorname{Pic}(X)\arrow[r]&\operatorname{Pic}(X_0)\arrow[r]&\operatorname{NS}(X_0).
\end{tikzcd}}
Soit enfin $d^t$ le pgcd des degrés des 0-cycles sur $X_\eta$. Avec ces notations, on a ce qui suit:

a) L'entier $d^t$ est aussi égal au pgcd des $\mu_i d_i$, où pour tout $i\in J$, $\mu_i$ désigne la multiplicité radicielle de $D_i$ sur $k$ (EGA IV 4.\quad); en particulier \deleted{$d$ divise $d^t$, donc est égal à 1 si $d^t$ l'est}\deleted{, \unsure{la réciproque étant vraie si $k$ est parfait}}\deleted{\inserted{et lui est égal si $k$ est parfait.}}\inserted{on a $d^t=p^h d$, où $p^h=\operatorname{pgcd}(\mu_i)$,}
\margin{\unsure{donc} $h=0$ i.e. $d=d^t$ si $k$ est parfait}
\margin{$\times$}

b) L'homomorphisme $w$ est un homomorphisme de $\underline Z$-modules \uline{de type fini dont le noyau est contenu dans le sous-groupe de torsion de $\Gamma$}. Ce dernier est isomorphe canoniquement à $\underline Z/d\underline Z$, avec comme générateur l'image dans $\Gamma$ du cycle $(1/d)X_0=\sum_{i\in J}(d_i/d)D_i$. Donc $\operatorname{Ker}w$ est cyclique

\margin{\figure{\begin{tikzpicture}\draw (0,8)--(0,0)--(.4,0);\draw[->](-.4,-.5)--(.4,0);\draw (0,-1)--(0,-2.3);\draw (-.7,8.3)--(-.4,8.9);\end{tikzpicture}}}
\mspage{60}
